% !PRINTED 556
\weberpage{556}
%<<CONT
taken absolutely, not greater than $k$. There is a number different from zero and divisible by $\mu$:

\weq{5}{\alpha = a_1 \omega_1 + a_2 \omega_2 + \cdots + a_n \omega_n,}

in which the integer coefficients $a_1, a_2 \ldots a_n$ do not exceed the limits $\pm k$, and consequently by (2) and (3)

\weq{6}{N_a(\alpha) < R k^n < R N_a(\mu).}

Now since $\alpha$ is divisible by $\mu$, we set

\weq{7}{\alpha = \mu\varphi,\ N_a(\alpha) = N_a(\mu) N_a(\varphi),}

and obtain from (6)

\weq{8}{N_a(\varphi) < R.}

If now $\psi$ is a representative of an arbitrarily given class $C$, we choose $\mu$ so that

\[
\beta = \mu\psi
\]

is a number, and if then, in accordance with (7),

\[
\alpha = \mu\varphi
\]

then

\[
\frac{\varphi}{\psi} = \frac{\alpha}{\beta},
\]

and thus $\varphi$ and $\psi$ are equivalent. Hence $\varphi$ is likewise a representative of the class $C$, and it satisfies the condition (8). Thus, as was to be proved, there occurs in every class a functional whose absolute norm lies below $R$.

The number of ideal classes, which we shall denote by $h$, is accordingly a natural number peculiar to the field $\Omega$, which is called the \gesperrt{class number}.

In the simplest case, where the class number is equal to 1, every functional is associated with a number; that is, every (integral or fractional) functional can be transformed into a unit by separating off a numerical factor. In this case every integer of the field $\Omega$ can be decomposed into \gesperrt{prime factors}, and these fields have a theory which agrees in essentials with rational number theory. For such fields the introduction of the functionals and ideals is not necessary.

Here belong, besides the field of the rational numbers, among others the field of the \gesperrt{Gauss} imaginary numbers and the field of the numbers formed from third roots of unity (Bd. I, \S.\,173, 174).
